\clearpage\thispagestyle{empty}
\pdfbookmark[0]{Sources for the Letter on Pallas and Juno}{pallasjuno18060708-sources}
\begin{center}\textbf{Sources for the Letter on Pallas and Juno}\end{center}
\noindent This AI translation follows Carl Friedrich Gauss’s letter of 8 July 1806 to Franz Xaver von Zach as printed in \href{https://gdz.sub.uni-goettingen.de/id/PPN236007467}{Gauss’s \textit{Werke}, vol.~VI} (Göttingen, 1874), pp.~277--278. It appeared in the \textit{Monatliche Correspondenz}, vol.~XIV, August 1806, article~XVI, part~1, \href{https://www.digitale-sammlungen.de/en/view/bsb10538606?page=201}{pp.~187}--\href{https://www.digitale-sammlungen.de/en/view/bsb10538606?page=206}{192}. The journal scans are from the Bayerische Staatsbibliothek.
\par\medskip
The collected edition gives a bracketed description of the Pallas ephemeris, which is printed in the journal on \href{https://www.digitale-sammlungen.de/en/view/bsb10538606?page=203}{pp.~189}--\href{https://www.digitale-sammlungen.de/en/view/bsb10538606?page=206}{192}. The cited Milan observations are on \href{https://www.digitale-sammlungen.de/en/view/bsb10538606?page=100}{p.~90 of the July issue}, where the journal credits Carlini with making and computing them. Friedrich Wilhelm Bessel’s letter and Juno observations begin on \href{https://www.digitale-sammlungen.de/en/view/bsb10538606?page=206}{p.~192 of the August issue}. Translation and typesetting produced by AI.
\clearpage
